% ===================================================================== PREAMBULE COMMUN — Carnet d'entraînement Fiche 9
\documentclass[11pt,a4paper]{article}
\usepackage[utf8]{inputenc}
\usepackage[T1]{fontenc}
\usepackage{lmodern}
\usepackage[french]{babel}
\usepackage{amsmath,amssymb,mathtools}
\usepackage{icomma}
\usepackage{siunitx}
\sisetup{output-decimal-marker={,},per-mode=symbol}
\DeclareSIUnit{\molar}{M}
\DeclareSIUnit{\Molar}{mol\,L^{-1}}
\usepackage[version=4]{mhchem}
\usepackage{xcolor}
\usepackage{colortbl}
\usepackage{tikz}
\usepackage{pgfplots}
\usepackage[most]{tcolorbox}
\usepackage{enumitem}
\usepackage{array}
\usepackage{booktabs}
\usepackage{multicol}
\usepackage{fancyhdr}
\usepackage[hidelinks]{hyperref}
\usetikzlibrary{arrows.meta,positioning,calc,decorations.pathmorphing,
  decorations.markings,angles,quotes,patterns,backgrounds,
  shapes.geometric,shapes.misc,fadings,intersections,fit}
\pgfplotsset{compat=1.18}
\usepackage[a4paper,margin=2.2cm,headheight=26pt,headsep=10pt]{geometry}

% ---------- Palette ----------
\definecolor{primary}{RGB}{150,98,162}
\definecolor{primarydark}{RGB}{92,52,110}
\definecolor{defcol}{RGB}{0,114,178}
\definecolor{thmcol}{RGB}{0,140,120}
\definecolor{formulacol}{RGB}{198,93,9}
\definecolor{remarkcol}{RGB}{120,94,200}
\definecolor{attncol}{RGB}{200,40,55}
\definecolor{excol}{RGB}{0,135,90}
\definecolor{methcol}{RGB}{90,90,100}
\definecolor{cationcol}{RGB}{0,90,190}
\definecolor{anioncol}{RGB}{200,60,55}
\definecolor{cristalcol}{RGB}{110,105,120}
\definecolor{eaucol}{RGB}{120,180,220}
\definecolor{solcol}{RGB}{0,135,90}
\definecolor{kpscol}{RGB}{198,93,9}
\definecolor{qcol}{RGB}{120,94,200}
\definecolor{precipcol}{RGB}{110,70,40}
\definecolor{exercol}{RGB}{150,98,162}
\definecolor{corrcol}{RGB}{0,120,100}

% ---------- Macros ----------
\newcommand{\Kps}{K_{\mathrm{ps}}}
\newcommand{\Ce}[1]{\left[\,\ce{#1}\,\right]}

% ---------- Style des boîtes ----------
\tcbset{boxbase/.style={enhanced,breakable,boxrule=0.9pt,arc=2.5pt,
   left=3mm,right=3mm,top=2.3mm,bottom=2.3mm,
   fonttitle=\bfseries,coltitle=white,toptitle=0.6mm,bottomtitle=0.6mm}}

\newtcolorbox[auto counter]{definition}[1][]{%
  boxbase,colback=defcol!5,colframe=defcol,
  title={\textsc{Définition}~\thetcbcounter\ \ifx\relax#1\relax\relax\else\textmd{--- \itshape #1}\fi}}
\newtcolorbox[auto counter]{theoreme}[1][]{%
  boxbase,colback=thmcol!6,colframe=thmcol,
  title={\textsc{Loi / Propriété}~\thetcbcounter\ \ifx\relax#1\relax\relax\else\textmd{--- \itshape #1}\fi}}
\newtcolorbox[auto counter]{formule}[1][]{%
  boxbase,colback=formulacol!6,colframe=formulacol,
  title={\textsc{Formule}~\thetcbcounter\ \ifx\relax#1\relax\relax\else\textmd{--- \itshape #1}\fi}}
\newtcolorbox{remarque}[1][]{%
  boxbase,colback=remarkcol!6,colframe=remarkcol,
  title={\textsc{Remarque}\ifx\relax#1\relax\relax\else\textmd{ : \itshape #1}\fi}}
\newtcolorbox{attention}[1][]{%
  boxbase,colback=attncol!6,colframe=attncol,
  title={\textsc{Attention}\ifx\relax#1\relax\relax\else\textmd{ : \itshape #1}\fi}}
\newtcolorbox{exemple}[1][]{%
  boxbase,colback=excol!5,colframe=excol,
  title={\textsc{Exemple}\ifx\relax#1\relax\relax\else\textmd{ : \itshape #1}\fi}}
\newtcolorbox{methode}[1][]{%
  boxbase,colback=methcol!7,colframe=methcol,sharp corners,
  title={\textsc{Méthode}\ifx\relax#1\relax\relax\else\textmd{ : \itshape #1}\fi}}
\newtcolorbox{astuce}[1][]{%
  boxbase,colback=primary!7,colframe=primary,
  title={\textsc{Astuce}\ifx\relax#1\relax\relax\else\textmd{ : \itshape #1}\fi}}

\newtcolorbox[auto counter]{exercice}[1][]{%
  boxbase,colback=exercol!4,colframe=exercol,
  title={\textsc{Exercice}~\thetcbcounter\ \ifx\relax#1\relax\relax\else\textmd{--- \itshape #1}\fi}}
\newtcolorbox[auto counter]{correction}[1][]{%
  boxbase,colback=corrcol!4,colframe=corrcol,
  title={\textsc{Corrigé}~\thetcbcounter\ \ifx\relax#1\relax\relax\else\textmd{--- \itshape #1}\fi}}

% ---------- Environnement « où : » ----------
\newenvironment{ou}{%
  \par\smallskip\noindent\textbf{où :}\nopagebreak
  \begin{itemize}[leftmargin=1.8em,topsep=2pt,itemsep=1.5pt,parsep=0pt]}%
  {\end{itemize}}

% ---------- Styles TikZ ----------
\tikzset{
  cat/.style={circle,fill=cationcol,draw=cationcol!50!black,inner sep=0pt,
              minimum size=5mm,font=\bfseries\scriptsize,text=white},
  an/.style={circle,fill=anioncol,draw=anioncol!50!black,inner sep=0pt,
             minimum size=5mm,font=\bfseries\scriptsize,text=white},
  crx/.style={circle,fill=cristalcol,draw=black!50,inner sep=0pt,
              minimum size=5mm,font=\bfseries\scriptsize,text=white},
  ptn/.style={circle,fill=black,inner sep=1.1pt}}

% ---------- En-têtes ----------
\pagestyle{fancy}
\fancyhf{}
\fancyhead[L]{\small\textcolor{primarydark}{\textbf{Fiche 9} — Solubilité et produit de solubilité}}
\fancyhead[R]{\small\textcolor{primarydark}{\textsc{Carnet d'entraînement}}}
\fancyfoot[C]{\small\thepage}
\renewcommand{\headrulewidth}{0.8pt}
\renewcommand{\headrule}{\hbox to\headwidth{\color{primary}\leaders\hrule height \headrulewidth\hfill}}
\usepackage{titlesec}
\titleformat{\section}{\Large\bfseries\color{primarydark}}{\thesection}{0.6em}{}
\titleformat{\subsection}{\large\bfseries\color{primary}}{\thesubsection}{0.6em}{}

% ---------- Bandeau de titre ----------
% #1 = thème/étiquette   #2 = titre principal   #3 = sous-titre
\newcommand{\bandeau}[3]{%
\begin{center}
\begin{tikzpicture}
\node[rectangle,rounded corners=7pt,fill=primary,text=white,
      minimum width=15.6cm,minimum height=1.95cm,align=center,inner sep=8pt]
  {{\large\bfseries #1}\\[3pt]{\Large\bfseries #2}\\[3pt]{\small #3}};
\end{tikzpicture}
\end{center}
\vspace{2mm}}
\begin{document}
\bandeau{Thème 3 — Quotient $Q$ \& prévision de précipitation}%
{Tuto — le critère $Q$ vs $\Kps$}%
{Décider si un mélange d'ions va précipiter, rester dissous, ou est saturé}

\noindent Le $\Kps$ décrit l'\emph{équilibre} (état saturé). Mais si l'on mélange des ions dans des
proportions \emph{quelconques}, va-t-il se former un précipité ? La réponse tient en \textbf{un seul
outil} : le quotient réactionnel $Q$, que l'on compare au $\Kps$.

\section*{1. Le quotient réactionnel $Q$}
\begin{definition}[quotient $Q$, ou produit ionique effectif]
$Q$ a la \textbf{même expression littérale} que le $\Kps$, mais avec les concentrations
\textbf{effectives} (celles réellement présentes dans le mélange, pas forcément à l'équilibre) :
\[ \boxed{\ Q = [\mathrm{M}^{q+}]_{\text{actuel}}^{\,p}\,[\mathrm{X}^{p-}]_{\text{actuel}}^{\,q}\ } \]
\end{definition}
\noindent \textbf{Différence clé} : $\Kps$ est une \emph{constante} (lue dans une table, fixée par $T$) ;
$Q$ est une \emph{variable} qui dépend des concentrations du moment.

\section*{2. Le critère de précipitation}
\begin{center}
\begin{tikzpicture}[scale=0.92]
  \draw[thick,->,>=Stealth] (-0.3,0) -- (12.6,0) node[right,font=\small]{$Q$};
  \fill[solcol!18] (0,0) rectangle (4.6,1.3);
  \node[solcol!40!black,font=\small\bfseries] at (2.3,0.85) {$Q<\Kps$};
  \node[solcol!40!black,font=\scriptsize] at (2.3,0.4) {non saturé : tout se dissout};
  \draw[thick,kpscol,line width=1.5pt] (4.6,0) -- (4.6,1.3);
  \node[kpscol,font=\small\bfseries] at (4.6,1.62) {$\Kps$};
  \node[kpscol!50!black,font=\scriptsize] at (4.6,-0.4) {saturé ($Q=\Kps$)};
  \fill[precipcol!22] (4.65,0) rectangle (12.2,1.3);
  \node[precipcol!40!black,font=\small\bfseries] at (8.4,0.85) {$Q>\Kps$};
  \node[precipcol!40!black,font=\scriptsize] at (8.4,0.4) {sursaturé : précipitation};
  \draw[->,>=Stealth,precipcol,thick] (8.4,-0.9) -- (5.2,-0.9)
       node[midway,below,precipcol!40!black,font=\scriptsize]{le solide se forme jusqu'à $Q=\Kps$};
\end{tikzpicture}
\end{center}

\begin{theoreme}[règle de comparaison]
\begin{itemize}[leftmargin=1.7em,topsep=1pt,itemsep=1pt]
  \item $Q<\Kps$ : solution \textbf{non saturée}, aucun précipité (un solide ajouté se dissoudrait).
  \item $Q=\Kps$ : solution \textbf{exactement saturée} (équilibre solide/solution).
  \item $Q>\Kps$ : solution \textbf{sursaturée}, il \textbf{se forme un précipité} jusqu'à ce que $Q$
        redescende à $\Kps$.
\end{itemize}
\end{theoreme}

\begin{methode}[décider s'il y a précipitation]
\textbf{(1)} repérer le sel peu soluble susceptible de se former, écrire son $\Kps$ ;
\textbf{(2)} calculer $Q$ avec les concentrations \emph{effectives} — \textbf{attention à la dilution}
si l'on mélange deux solutions ! ; \textbf{(3)} comparer $Q$ et $\Kps$.
\end{methode}

\begin{attention}[la dilution lors d'un mélange]
Quand on mélange un volume $V_1$ d'une solution avec un volume $V_2$ d'une autre, chaque concentration
est \textbf{recalculée} sur le volume total $V_1+V_2$ :
\[ [\text{ion}]_{\text{après}} = \dfrac{C_{\text{avant}}\times V_{\text{de cet ion}}}{V_1+V_2}. \]
Oublier cette dilution est l'erreur la plus fréquente des exercices de précipitation.
\end{attention}

\begin{exemple}[\ce{Mg(OH)2} précipite-t-il dans l'eau de mer ?]
Données : $[\ce{Mg^2+}]=\qty{0{,}06}{\molar}$, $[\ce{OH-}]=\qty{2{,}0e-3}{\molar}$,
$\Kps(\ce{Mg(OH)2})=\num{9{,}0e-12}$.
\[ Q=[\ce{Mg^2+}][\ce{OH-}]^{2}=(0{,}06)(2{,}0\times10^{-3})^{2}=\num{2{,}4e-7}. \]
$Q=\num{2{,}4e-7}\gg\Kps=\num{9{,}0e-12}$ : la solution est largement sursaturée, \ce{Mg(OH)2}
\textbf{précipite} (procédé d'extraction du magnésium marin).
\end{exemple}

\end{document}
